
\section*{CHƯƠNG 2 - CƠ SỞ LÝ THUYẾT}
\addcontentsline{toc}{section}{\numberline{}CHƯƠNG 2 - CƠ SỞ LÝ THUYẾT}
\setcounter{section}{2}
\setcounter{subsection}{0}
\setcounter{figure}{0}
\setcounter{table}{0}

\subsection{Tổng quan về Sentiment Analysis}

\subsubsection{Khái niệm Sentiment Analysis}
Sentiment Analysis (Phân tích cảm xúc) là một lĩnh vực của xử lý ngôn ngữ tự nhiên (Natural Language Processing - NLP) và khai phá văn bản (Text Mining), tập trung vào việc xác định, trích xuất và phân loại ý kiến, cảm xúc, thái độ của người viết trong văn bản. Sentiment Analysis còn được gọi là Opinion Mining, cho phép máy tính hiểu được cảm xúc tích cực, tiêu cực hay trung lập trong các đoạn văn bản.

Trong bối cảnh hiện đại, với sự bùng nổ của mạng xã hội và các nền tảng trực tuyến, Sentiment Analysis đã trở thành công cụ quan trọng giúp doanh nghiệp hiểu rõ ý kiến khách hàng, theo dõi thương hiệu và đưa ra quyết định kinh doanh dựa trên dữ liệu.

\subsubsection{Các mức độ phân tích cảm xúc}
Sentiment Analysis có thể được thực hiện ở nhiều mức độ khác nhau:

\textbf{Document-level:} Phân tích cảm xúc của toàn bộ văn bản (document), xác định văn bản đó mang cảm xúc tích cực hay tiêu cực tổng thể.

\textbf{Sentence-level:} Phân tích cảm xúc của từng câu riêng lẻ trong văn bản.

\textbf{Aspect-level:} Phân tích cảm xúc đối với các khía cạnh cụ thể của đối tượng được đề cập (ví dụ: chất lượng phim, diễn xuất, kịch bản).

\textbf{Entity-level:} Phân tích cảm xúc đối với các thực thể cụ thể được nhắc đến trong văn bản.

Đề tài này tập trung vào Document-level Sentiment Analysis, phân loại toàn bộ bình luận thành ba nhãn: Positive, Negative và Neutral.

\subsubsection{Ứng dụng của Sentiment Analysis}
\begin{itemize}
    \item \textbf{Marketing và Thương hiệu:} Theo dõi phản ứng khách hàng về sản phẩm, dịch vụ, chiến dịch quảng cáo.
    \item \textbf{Dịch vụ khách hàng:} Phát hiện khách hàng không hài lòng để xử lý kịp thời.
    \item \textbf{Phân tích thị trường:} Dự đoán xu hướng thị trường dựa trên sentiment của người tiêu dùng.
    \item \textbf{Chính trị:} Phân tích ý kiến công chúng về chính sách, ứng cử viên.
    \item \textbf{Giải trí:} Đánh giá phản ứng khán giả về phim, nhạc, game.
\end{itemize}

\subsection{Các phương pháp Sentiment Analysis}

\subsubsection{Rule-based Approach}
Phương pháp dựa trên quy tắc sử dụng các từ điển cảm xúc (sentiment lexicon) được định nghĩa trước. Mỗi từ trong từ điển được gán một điểm cảm xúc (sentiment score).

\textbf{Ưu điểm:}
\begin{itemize}
    \item Đơn giản, dễ hiểu và triển khai
    \item Không cần dữ liệu huấn luyện
    \item Giải thích được kết quả
\end{itemize}

\textbf{Nhược điểm:}
\begin{itemize}
    \item Không xử lý được ngữ cảnh phức tạp
    \item Khó xử lý từ lóng, từ mới
    \item Độ chính xác hạn chế
\end{itemize}

\subsubsection{Machine Learning Approach}
Phương pháp học máy sử dụng các thuật toán như Naive Bayes, SVM, Random Forest để học từ dữ liệu đã được gán nhãn.

\textbf{Ưu điểm:}
\begin{itemize}
    \item Độ chính xác cao hơn rule-based
    \item Tự động học từ dữ liệu
    \item Xử lý được ngữ cảnh tốt hơn
\end{itemize}

\textbf{Nhược điểm:}
\begin{itemize}
    \item Cần dữ liệu huấn luyện lớn
    \item Phụ thuộc vào chất lượng dữ liệu
    \item Khó giải thích kết quả
\end{itemize}

\subsubsection{Deep Learning Approach}
Phương pháp học sâu sử dụng các mô hình như LSTM, BERT, RoBERTa để phân tích cảm xúc với độ chính xác cao.

\textbf{Ưu điểm:}
\begin{itemize}
    \item Độ chính xác rất cao
    \item Xử lý ngữ cảnh phức tạp
    \item Hiểu được ý nghĩa sâu của văn bản
\end{itemize}

\textbf{Nhược điểm:}
\begin{itemize}
    \item Cần tài nguyên tính toán lớn
    \item Thời gian huấn luyện lâu
    \item Cần dữ liệu huấn luyện rất lớn
\end{itemize}

\subsection{YouTube Data API v3}

\subsubsection{Giới thiệu YouTube Data API}
YouTube Data API v3 là giao diện lập trình ứng dụng (API) do Google cung cấp, cho phép các nhà phát triển truy cập và thao tác với dữ liệu YouTube một cách có cấu trúc. API này cung cấp khả năng tìm kiếm video, lấy thông tin video, quản lý playlist, và đặc biệt là thu thập bình luận.

\subsubsection{Các tính năng chính}
\begin{itemize}
    \item \textbf{Search:} Tìm kiếm video, channel, playlist theo từ khóa
    \item \textbf{Videos:} Lấy thông tin chi tiết về video (tiêu đề, mô tả, thống kê)
    \item \textbf{Comments:} Thu thập bình luận và phản hồi
    \item \textbf{Channels:} Lấy thông tin về kênh YouTube
    \item \textbf{Playlists:} Quản lý danh sách phát
\end{itemize}

\subsubsection{Cơ chế hoạt động}
YouTube Data API sử dụng giao thức RESTful HTTP, cho phép gửi request và nhận response dưới dạng JSON. Mỗi request cần có API key để xác thực.

\textbf{Quota System:}
YouTube API có hệ thống quota để giới hạn số lượng request. Mỗi API key miễn phí có 10,000 units/ngày. Các thao tác khác nhau tiêu tốn số units khác nhau:
\begin{itemize}
    \item Read operation: 1 unit
    \item Write operation: 50 units
    \item Video upload: 1600 units
\end{itemize}

\textbf{Pagination:}
Khi dữ liệu trả về lớn, API sử dụng phân trang (pagination) với tham số \texttt{pageToken} để lấy trang tiếp theo.

\subsubsection{Thu thập bình luận}
Để thu thập bình luận, sử dụng endpoint \texttt{commentThreads.list} với các tham số:
\begin{itemize}
    \item \texttt{part}: Phần dữ liệu cần lấy (snippet, replies)
    \item \texttt{videoId}: ID của video
    \item \texttt{maxResults}: Số lượng kết quả tối đa (1-100)
    \item \texttt{pageToken}: Token để lấy trang tiếp theo
    \item \texttt{textFormat}: Định dạng văn bản (plainText hoặc html)
\end{itemize}

\subsection{Xử lý ngôn ngữ tự nhiên (NLP)}

\subsubsection{Khái niệm NLP}
Natural Language Processing (NLP) là lĩnh vực nghiên cứu về cách máy tính có thể hiểu, phân tích và tạo ra ngôn ngữ tự nhiên của con người. NLP kết hợp kiến thức từ ngôn ngữ học, khoa học máy tính và trí tuệ nhân tạo.

\subsubsection{Các kỹ thuật tiền xử lý văn bản}
Tiền xử lý văn bản là bước quan trọng trong NLP, giúp làm sạch và chuẩn hóa dữ liệu trước khi phân tích.

\textbf{Tokenization:} Chia văn bản thành các đơn vị nhỏ hơn (tokens) như từ, câu.

\textbf{Lowercasing:} Chuyển tất cả ký tự về chữ thường để đảm bảo tính nhất quán.

\textbf{Removing Special Characters:} Loại bỏ các ký tự đặc biệt, số, dấu câu không cần thiết.

\textbf{Removing URLs and Mentions:} Loại bỏ đường link và mention (@username) không mang ý nghĩa cảm xúc.

\textbf{Removing Stop Words:} Loại bỏ các từ phổ biến không mang nhiều ý nghĩa (the, is, at, which).

\textbf{Stemming/Lemmatization:} Đưa từ về dạng gốc (running → run).

\textbf{Removing Duplicates:} Loại bỏ các văn bản trùng lặp.

\subsection{TextBlob - Thư viện phân tích cảm xúc}

\subsubsection{Giới thiệu TextBlob}
TextBlob là một thư viện Python đơn giản để xử lý dữ liệu văn bản. Nó cung cấp API đơn giản cho các tác vụ NLP phổ biến như part-of-speech tagging, noun phrase extraction, sentiment analysis, classification, translation.

\subsubsection{Sentiment Analysis với TextBlob}
TextBlob sử dụng phương pháp rule-based kết hợp với lexicon để phân tích cảm xúc. Nó trả về hai giá trị:

\textbf{Polarity (Độ cực):} Giá trị từ -1.0 đến 1.0
\begin{itemize}
    \item -1.0: Rất tiêu cực
    \item 0.0: Trung lập
    \item 1.0: Rất tích cực
\end{itemize}

\textbf{Subjectivity (Độ chủ quan):} Giá trị từ 0.0 đến 1.0
\begin{itemize}
    \item 0.0: Rất khách quan (fact-based)
    \item 1.0: Rất chủ quan (opinion-based)
\end{itemize}

\subsubsection{Cách hoạt động}
TextBlob sử dụng từ điển cảm xúc được xây dựng sẵn, trong đó mỗi từ được gán một điểm polarity và subjectivity. Khi phân tích một câu, TextBlob:
\begin{enumerate}
    \item Tokenize câu thành các từ
    \item Tra cứu điểm cảm xúc của mỗi từ trong từ điển
    \item Xử lý các modifier (very, not, extremely) để điều chỉnh điểm
    \item Tính trung bình điểm của tất cả các từ
    \item Trả về polarity và subjectivity cuối cùng
\end{enumerate}

\subsection{Metrics đánh giá}

\subsubsection{Confusion Matrix}
Confusion Matrix là bảng ma trận thể hiện số lượng dự đoán đúng và sai của mô hình phân loại.

\begin{table}[H]
\centering
\caption{Confusion Matrix cho phân loại 3 nhãn}
\begin{tabular}{|c|c|c|c|}
\hline
\textbf{Actual / Predicted} & \textbf{Positive} & \textbf{Negative} & \textbf{Neutral} \\
\hline
\textbf{Positive} & TP & FN & FN \\
\hline
\textbf{Negative} & FP & TN & FN \\
\hline
\textbf{Neutral} & FP & FP & TN \\
\hline
\end{tabular}
\end{table}

\subsubsection{Accuracy}
Accuracy (Độ chính xác) là tỷ lệ dự đoán đúng trên tổng số dự đoán:

\[
\text{Accuracy} = \frac{\text{Số dự đoán đúng}}{\text{Tổng số dự đoán}}
\]

\subsubsection{Precision và Recall}
\textbf{Precision (Độ chính xác):} Tỷ lệ dự đoán đúng trong số các dự đoán positive:

\[
\text{Precision} = \frac{TP}{TP + FP}
\]

\textbf{Recall (Độ phủ):} Tỷ lệ dự đoán đúng trong số các mẫu thực sự positive:

\[
\text{Recall} = \frac{TP}{TP + FN}
\]

\subsubsection{F1-Score}
F1-Score là trung bình điều hòa của Precision và Recall:

\[
\text{F1-Score} = 2 \times \frac{\text{Precision} \times \text{Recall}}{\text{Precision} + \text{Recall}}
\]

\subsection{Trực quan hóa dữ liệu}

\subsubsection{Matplotlib}
Matplotlib là thư viện Python phổ biến nhất để tạo biểu đồ và trực quan hóa dữ liệu. Nó cung cấp API tương tự MATLAB, cho phép tạo các loại biểu đồ: line plot, bar chart, scatter plot, histogram, pie chart.

\subsubsection{Seaborn}
Seaborn là thư viện trực quan hóa dữ liệu được xây dựng trên Matplotlib, cung cấp giao diện cấp cao hơn và các biểu đồ thống kê đẹp mắt hơn. Seaborn đặc biệt hữu ích cho:
\begin{itemize}
    \item Biểu đồ phân bố (distribution plots)
    \item Biểu đồ tương quan (correlation plots)
    \item Biểu đồ categorical (box plot, violin plot)
    \item Heatmap
\end{itemize}

\subsection{Tổng kết}
Chương này đã trình bày các kiến thức nền tảng cần thiết cho đề tài, bao gồm:
\begin{itemize}
    \item Khái niệm và các phương pháp Sentiment Analysis
    \item YouTube Data API v3 và cách thu thập dữ liệu
    \item Các kỹ thuật tiền xử lý văn bản trong NLP
    \item TextBlob và cách phân tích cảm xúc
    \item Các metrics đánh giá hiệu năng mô hình
    \item Công cụ trực quan hóa dữ liệu
\end{itemize}

Những kiến thức này sẽ là cơ sở để thiết kế và triển khai hệ thống trong các chương tiếp theo.
